\chapter{References and divergences}
\label{chap:references}

This chapter names every published source lean4lean's kernel layer is checked against, places
lean4lean on the same grid as the Rocq development it parallels, and lists, notion by notion,
every point where MetaRocq/PCUIC, Carneiro's thesis and lean4lean's kernel diverge.

\section{The sources}
\label{sec:ref-sources}

\subsection{Carneiro's thesis [CT]}

\begin{itemize}
\item \lead{Identity} \emph{The Type Theory of Lean}, Mario Carneiro, MSc draft, Carnegie Mellon
University; a \LaTeX{} source tree (\texttt{axioms.tex}, \texttt{typesys.tex},
\texttt{unique.tex}, \texttt{Wtypes.tex}, \texttt{soundness.tex}, \texttt{normalization.tex},
\texttt{compilation.tex}), not a rendered, finished document.
\item \lead{Establishes} \texttt{axioms.tex}: typing, definitional equality $\equiv$, algorithmic
equality $\Leftrightarrow$, reduction, universes with \texttt{imax}, let/delta reduction, the
inductive \texttt{spec}/\texttt{ctor} judgments and the recursor's type, $\iota$ and K-like
reduction, quotients, propositional extensionality, choice.
\item \lead{Establishes} \texttt{typesys.tex}: undecidability of $\equiv$, non-transitivity of
$\Leftrightarrow$, failure of subject reduction for the algorithmic judgment, plus
Regularity/Weakening/Substitution.
\item \lead{Establishes} \texttt{unique.tex}: unique typing (\texttt{thm:unique}) via an
$n$-indexed joint induction with Church-Rosser for a restricted $\kappa$ reduction
(\texttt{thm:church\_rosser}), the completeness of that factoring (\texttt{thm:ckappa}), and
definitional inversion (\texttt{thm:1dinv}).
\item \lead{Establishes} \texttt{Wtypes.tex}: a $\Sigma$-projection typing shape for structures,
with an $\eta$ rule the text itself flags as not a real Lean definitional equality.
\item \lead{Cited by lean4lean} the K-A rows of \S\ref{sec:ref-divergences} cite
\texttt{axioms.tex} throughout; \texttt{typesys.tex}'s three lemmas ground
Weakening/Substitution/Regularity; \texttt{unique.tex}'s \texttt{thm:unique} and
\texttt{thm:church\_rosser} are the targets of K-C3 and K-C1.
\item \lead{Unfinished} \texttt{normalization.tex} ends mid-section, literal token
\texttt{UNFINISHED}, after introducing only the restricted $K$ (not $K^+$) reduction, with no
termination theorem. \texttt{soundness.tex}'s reflection extension stops mid-sentence, same
token. \texttt{compilation.tex} is a 25-line fragment with no accompanying theorem.
\end{itemize}

\subsection{The MetaRocq checker verification: [JACM] and [CCC]}

\begin{itemize}
\item \lead{Identity} [CCC]: Sozeau, Boulier, Forster, Tabareau, Winterhalter, "Coq Coq
Correct!", POPL 2020, PACMPL 4, Article~8, 28~pp. [JACM]: Sozeau, Forster, Lennon-Bertrand,
Nielsen, Tabareau, Winterhalter, J.~ACM 72(1), Article~8, 74~pp, Jan 2025.
\item \lead{Relation} [JACM] p.~8:5, "Prior publication," names [CCC] as the conference version
it extends: [CCC] is \hl{earlier}, not later, and its Sec 2-3 already cover PCUIC and kernel
verification in full.
\item \lead{Establishes -- kernel theory} [CCC] Sec 2 introduces PCUIC: syntax, universes and
valuations, cumulativity, typing, confluence by the triangle method. Subject reduction,
validity, principal typing, strong normalisation and the guard condition are
\texttt{Conjecture}/\texttt{Axiom}, not proved.
\item \lead{Establishes -- kernel verification} [CCC] Sec 3: a fuel-free weak-head reduction
machine, a graph-based universe-consistency decision procedure, a conversion checker, and a
bidirectional \texttt{infer} algorithm, all proved sound; completeness is not claimed.
\item \lead{Establishes -- kernel theory} [JACM] Sec 3-4: the same PCUIC specification plus an
algorithmic specification (acyclic-graph universes, reduction-based conversion, bidirectional
typing) proved equivalent to it, and the first treatment of cumulative inductive types with a
subject-reduction proof.
\item \lead{Establishes -- kernel verification} [JACM] Sec 6, pp.~8:56-8:59: the checker is
proved \hl{sound and complete} against the declarative specification, every rejection carrying a
refutation of the corresponding derivation; the method found and fixed a genuine completeness
bug in Coq 8.14.
\item \lead{Establishes -- metatheory} [JACM] Sec 5: confluence, subject reduction, weak-CBV
standardisation (5.6, the one use of the strong-normalisation axiom in this part of the
development), canonicity (5.7), consistency (5.8).
\item \lead{Cited by lean4lean} every K-A/K-B/K-C/K-D/K-E row of \S\ref{sec:ref-divergences}
cites [JACM] Sec 3-6 as the default reference; [CCC] only where its earlier framing (Sec 2-3,
Fig.~1-10) is itself the artefact under comparison.
\end{itemize}

\subsection{The MetaRocq project overview [MR]}

\begin{itemize}
\item \lead{Identity} Sozeau, Anand, Boulier, Cohen, Forster, Kunze, Malecha, Tabareau,
Winterhalter, "The MetaCoq Project", J.~Automated Reasoning 64:947-999, 2020; the extended
journal version of an ITP 2018 paper.
\item \lead{Not a verification result} Sect.~2.7, p.~973 states explicitly that its
\texttt{infer}/\texttt{check}/\texttt{check\_conv} algorithms have \hl{no soundness or
completeness proof}: this is the pre-verification stage [CCC] and [JACM] later prove correct.
\item \lead{Establishes} Sect.~2: a reification \texttt{term} of Coq's kernel syntax (Fig.~1),
its typing judgment, reduction, well-formed environments and universe algebra -- the
pre-PCUIC ancestor of [CCC]'s specification. The string "PCUIC" never occurs in this paper.
\item \lead{Cited by lean4lean} K-E2's evidence line cites [MR] Sect.~3, pp.~973-977, for the
unverified Template quoting layer, and Sect.~2.7, p.~973, for the "no proof yet" baseline the
later checkers supersede.
\end{itemize}

\subsection{The habilitation [HAB]}

\begin{itemize}
\item \lead{Identity} Matthieu Sozeau, \emph{Mechanized Metatheory and Certified Compilation for
The Rocq Prover}, HDR, Universite de Nantes, 2026, HAL tel-05544162, 152~pp; a retrospective
synthesis manuscript, not a fresh research paper.
\item \lead{Superset} lists [JACM] first among four articles it "is based on"; Ch.~2-5 are the
current-naming (MetaRocq v1.4, Rocq 9.0) rewrite of [JACM]/[CCC]'s kernel theory and
verification.
\item \lead{Establishes -- kernel theory} Ch.~2, pp.~26-42 (Specs 1-2, 18-19): PCUIC syntax,
universes, cumulativity, typing, well-formed global environments including positivity and
cumulative inductive types.
\item \lead{Establishes -- kernel verification} Ch.~3-5, pp.~44-84 (Specs 4, 6, 7, 8, 16-17):
the algorithmic specification and its equivalence to the declarative one, core metatheory, and
the sound-and-complete checker.
\item \lead{Cited by lean4lean} the fallback citation throughout \S\ref{sec:ref-divergences} for
current MetaRocq naming, e.g.\ K-A1's universe algebra, K-A5's eta rule, K-D1's checker
guarantee.
\item \lead{Where thin} Ch.~1 is an informal tour with no theorem; Ch.~8's open-problems list,
pp.~127-131, names eta, \texttt{SProp} extensionality and guard-checker verification as unproved.
\end{itemize}

\subsection{The MetaRocq docs [MRD]}

\begin{itemize}
\item \lead{Identity} \texttt{metarocq.github.io}, the project's live documentation: a landing
page plus per-package \texttt{README.md} files and a generated coqdoc index; no version number,
and it drifts ahead of every paper.
\item \lead{Establishes} the file/module map for the kernel layer: \texttt{PCUICAst},
\texttt{PCUICTyping}, \texttt{PCUICCumulativitySpec}, \texttt{PCUICConfluence}, \texttt{PCUICSR},
\texttt{PCUICValidity}, \texttt{PCUICPrincipality}, \texttt{PCUICSN} (\texttt{pcuic/theories/}),
and \texttt{PCUICSafeReduce}, \texttt{PCUICSafeConversion}, \texttt{PCUICSafeRetyping},
\texttt{PCUICSafeChecker} (\texttt{safechecker/theories/}).
\item \lead{Cited by lean4lean} the tiebreaker whenever \S\ref{sec:ref-divergences} needs a
current name or line number, e.g.\ \texttt{check\_wf\_env} at
\texttt{PCUICSafeChecker.v:2398} in K-A12/K-E3.
\item \lead{Caveat} no theorem numbers and no page numbers, only a one-line-per-file prose
index; always cited by file name, never by figure or theorem number.
\end{itemize}

\section{Where lean4lean sits}
\label{sec:ref-where}

Restricted to the kernel layer, the project grid has two columns and two rows.

\par\smallskip\noindent
\begin{tabular}{lp{0.43\linewidth}p{0.43\linewidth}}
\toprule
 & Kernel theory & Kernel verification \\
\midrule
Rocq & PCUIC's declarative specification: universes and valuations, cumulativity, typing,
well-formed global environments ([JACM] Sec 3-4; [HAB] Ch.~2). &
MetaRocq's SafeChecker: fuel-free weak-head reduction, a universe-consistency decision
procedure, a checker proved sound \emph{and} complete ([JACM] Sec 6; [HAB] Ch.~3-5). \\
Lean & Carneiro's thesis [CT]: typing, $\equiv$, $\Leftrightarrow$, reduction, universes with
\texttt{imax}, the inductive/recursor calculus, unique typing, Church-Rosser for $\kappa$. &
lean4lean's executable kernel, its correctness stated relative to
\texttt{VExpr}/\texttt{VEnv}/\texttt{IsDefEq}/\texttt{HasType} (\ref{def:isdefeq}) through a
translation \texttt{TrExprS} (K-D2), sound but not complete. \\
\bottomrule
\end{tabular}\par\smallskip

\begin{itemize}
\item [MR]: Rocq kernel theory only (the pre-PCUIC \texttt{term}), explicitly not a verification
result.
\item [CCC]: first Rocq kernel verification (sound only), plus the introduction of PCUIC itself.
\item [JACM]: Rocq kernel theory and verification, sound and complete, conjecture-free.
\item [HAB]: a superset of all three, in current MetaRocq naming.
\item [MRD]: the file-and-module map, the tiebreaker between a paper and the current code.
\item [CT]: the Lean kernel theory in full; no Lean kernel-verification result of its own, only
negative bounds on the algorithmic judgment.
\end{itemize}
\paragraph{Verification method.} \textbf{MetaRocq} writes its checker in Rocq against the same
syntax its specification uses, proves it sound and complete, then extracts it: the checked
object and the specified object are the same terms, with no refinement layer (K-D2).
\textbf{lean4lean} ports Lean's C++ kernel onto the real \texttt{Lean.Expr}/
\texttt{Lean.Kernel.Environment} and proves soundness relative to the model through a
translation \texttt{TrExprS}/\texttt{TrEnv} (K-D2) -- a refinement layer with its own trust
boundary (K-D4, K-D5). It runs on explicit fuel counters rather than a normalisation axiom
(K-C5, K-D7), and claims soundness only: [CT]'s \texttt{typesys.tex} proves $\equiv$
undecidable and $\Leftrightarrow$ non-transitive, so no algorithm can decide either (K-D1).

\section{Three-way divergences}
\label{sec:ref-divergences}

Fifty-two items compare MetaRocq/PCUIC, Carneiro's thesis and lean4lean's kernel layer, notion
by notion; each is read off a source file, a blueprint label, or the project's own status
chapter, never inferred from a paper alone. Five kinds, with their count among the 52:

\begin{itemize}
\item \textbf{FORCED-LEAN} (7): Lean's type theory itself has or lacks the feature; no design
choice is available.
\item \textbf{FORCED-LEAN4LEAN} (2): the port-and-refine verification method, not Lean's theory,
forces the divergence.
\item \textbf{DESIGN} (22): a deliberate modelling choice, made where either an equivalent or a
stronger route was open.
\item \textbf{GAP} (19): an open proof obligation, a missing derivation, or a statement not yet
proved.
\item \textbf{OPEN} (2): out of scope by explicit statement, with no route to close it planned.
\end{itemize}

\subsection{K-A: the object language}

\par\smallskip\noindent
\begin{tabular}{lp{0.13\linewidth}p{0.29\linewidth}p{0.26\linewidth}l}
\toprule
id & notion & PCUIC/thesis & lean4lean & kind \\
\midrule
K-A1 & universes & sorts over a constraint store; thesis: levels $0,S,\max,\mathrm{imax}$, no
\texttt{SProp} & \texttt{VLevel}: zero/succ/max/imax/param, one sort family & FORCED-LEAN \\
K-A2 & level order & graph procedure over valuations; thesis: nine syntactic rules for
$\ell\le\ell'+n$ & semantic \texttt{VLevel.LE} over all valuations & DESIGN \\
K-A3 & cumulativity & directed \texttt{cumulSpec0}; thesis: no subtyping rule at all &
\texttt{IsDefEq} is symmetric, no cumulativity & FORCED-LEAN \\
K-A4 & proof irrelevance & no rule (only in the checker's guard); thesis: any two proofs of
\texttt{p} are $\equiv$ & \texttt{IsDefEq.proofIrrel}, the same rule & FORCED-LEAN \\
K-A5 & eta, functions & out of scope for PCUIC; thesis: $e:\forall y{:}\alpha.\beta \Rightarrow
\lambda x.(ex)\equiv e$ & \texttt{IsDefEq.eta}, the same rule & FORCED-LEAN \\
K-A6 & eta, structures & no rule, no obligation on either reference side & the checker has it,
the model does not; both \texttt{sorry} & GAP \\
K-A7 & let/zeta & \texttt{tLetIn} with a zeta rule; thesis: zeta kept, with a conservativity
argument & no let in \texttt{VExpr}; substituted away at translation time & DESIGN \\
K-A8 & projections & \texttt{tProj} with its own typing/iota/congruence rules; thesis: a
Sigma-projection shape, no eta & \texttt{TrProj}: a recursor application via the iota registry &
DESIGN \\
K-A9 & projection coverage & every declared projection typechecks; no thesis counterpart &
\texttt{TrProj} covers only non-mutual, single-constructor shapes; \texttt{sorry} & GAP \\
K-A10 & iota, K-like reduction & one \texttt{cumul\_iota} rule, no K; thesis: constructor rule
plus K-like for subsingletons & one \texttt{pats} registry rule; K-like omitted & GAP \\
K-A11 & quotients & none in PCUIC; thesis: \texttt{Quot}/\texttt{mk}/\texttt{sound}/\texttt{lift},
"semi-builtin" & \texttt{quotDefEq}, one \texttt{VDefEq} axiom, not a \texttt{pats} entry &
DESIGN \\
K-A12 & inductive well-formedness & established by \texttt{check\_wf\_env}; thesis: generated
from the block spec & \texttt{VInductDecl.WF}, 20 fields, asserted not derived & GAP \\
K-A13 & primitives & \texttt{tPrim} with its own typing rule; no thesis counterpart & no literal
node; a literal unfolds to constructors & DESIGN \\
K-A14 & native reduction & not modelled on either reference side & \texttt{reduceNative} refuses
\texttt{reduceBool}; vacuous & OPEN \\
\bottomrule
\end{tabular}\par\smallskip
\paragraph{K-A1: universes and their algebra}
\begin{itemize}
\item \lead{PCUIC} \texttt{sort} $=\;$\texttt{sProp}$\,|\,$\texttt{sSProp}$\,|\,$
\texttt{sType\ univ}, a set of level/sign pairs read through a global valuation and constraint
store.
\item \lead{Thesis} levels $\ell ::= u \mid 0 \mid S\ell \mid \max(\ell,\ell) \mid
\mathrm{imax}(\ell,\ell)$, with a $\mathbb{N}$-valued semantics and no \texttt{SProp}.
\item \lead{lean4lean} \texttt{VLevel} (\ref{ind:vlevel}): zero, succ, max, imax, \texttt{param}
as a de~Bruijn index; one sort family \texttt{VExpr.sort}.
\item \lead{Kind} FORCED-LEAN.
\item \lead{Why} Lean's universes are closed level expressions over a declaration's own
parameters, with no constraint graph and no \texttt{Set}/\texttt{SProp}.
\item \lead{Consequence} Judgments are indexed by a universe arity, not a constraint set, and
polymorphism is level instantiation rather than \texttt{consistent\_instance\_ext}.
\end{itemize}
\paragraph{K-A2: the level order is semantic, not a rule system}
\begin{itemize}
\item \lead{PCUIC} \texttt{leq\_universe} over valuations plus a graph procedure for constraint
consistency.
\item \lead{Thesis} $\ell\le\ell'+n$ by nine syntactic rules, including three \texttt{imax}
identities and a case split on a parameter.
\item \lead{lean4lean} \texttt{VLevel.LE}$\,a\,b := \forall\,ls,\ a.\mathrm{eval}\,ls \le
b.\mathrm{eval}\,ls$ (\ref{def:vlevel-le-equiv}).
\item \lead{Kind} DESIGN.
\item \lead{Why} a semantic definition needs no rule-completeness proof; the thesis's
\texttt{imax} identities become one \texttt{simp} lemma each.
\item \lead{Consequence} two sorts are equal exactly when they denote the same function of the
parameters, so the decision problem is genuinely hard (K-T2).
\end{itemize}
\paragraph{K-A3: no cumulativity}
\begin{itemize}
\item \lead{PCUIC} directed \texttt{cumulSpec0}, with \texttt{cumul\_Sort}, \texttt{cumul\_Ind},
\texttt{cumul\_Construct} for cumulative inductive types.
\item \lead{Thesis} no directed relation at all; equality is the only comparison stated.
\item \lead{lean4lean} \texttt{IsDefEq} is an equivalence; \texttt{defeqDF} transports along a
definitional equality of types, not a lower bound.
\item \lead{Kind} FORCED-LEAN.
\item \lead{Why} Lean has neither universe cumulativity nor cumulative inductive types.
\item \lead{Consequence} unique typing gives an equality of the two types directly, needing no
"unique sort quality" strengthening.
\end{itemize}
\paragraph{K-A4: definitional proof irrelevance}
\begin{itemize}
\item \lead{PCUIC} no rule; \texttt{sSProp} exists as a sort but conversion never identifies its
inhabitants.
\item \lead{Thesis} from $p{:}P$, $h{:}p$, $h'{:}p$, derive $h\equiv h'$, for any \texttt{Prop}.
\item \lead{lean4lean} \texttt{IsDefEq.proofIrrel}: the same rule, reading the shared type of
both sides.
\item \lead{Kind} FORCED-LEAN.
\item \lead{Why} Lean's kernel identifies any two proofs of a \texttt{Prop}.
\item \lead{Consequence} conversion must be typed (K-B2); ordinary Church-Rosser needs a
separate \texttt{NormalEq} congruence absorbing irrelevance and eta.
\end{itemize}
\paragraph{K-A5: eta for functions as a primitive rule}
\begin{itemize}
\item \lead{PCUIC} no eta rule; explicitly named as out of scope.
\item \lead{Thesis} $e:\forall y{:}\alpha.\beta \Rightarrow \lambda x{:}\alpha.(e\,x)\equiv e$.
\item \lead{lean4lean} \texttt{IsDefEq.eta}, the same rule, read against $e$'s
\texttt{forallE} type.
\item \lead{Kind} FORCED-LEAN.
\item \lead{Why} Lean's kernel has definitional eta for functions.
\item \lead{Consequence} eta cannot be a reduction; it enters \texttt{NormalEq} as two directed
rules, needing its own symmetry theorem and an eta-expansion-aware confluence argument.
\end{itemize}
\paragraph{K-A6: structure eta and unit-like types are not modelled}
\begin{itemize}
\item \lead{PCUIC} neither structure eta nor a unit-like rule exists; the two reference sides
agree, so no obligation arises.
\item \lead{Thesis} no counterpart.
\item \lead{lean4lean} the executable kernel implements both (\ref{thm:vtc-tryetastruct},
\ref{thm:vtc-isdefequnitlike}); \texttt{IsDefEq} has neither rule, and both obligations are
literal \texttt{sorry}.
\item \lead{Kind} GAP.
\item \lead{Why} the statement is not derivable from the current rule set; a proof-effort gap
alone would not close it.
\item \lead{Consequence} the checker's three top-level entry points are partial as stated: the
checker can accept pairs the model cannot relate.
\end{itemize}
\paragraph{K-A7: let is substituted away; no zeta rule}
\begin{itemize}
\item \lead{PCUIC} \texttt{tLetIn} as a term former, with both a zeta reduction and a
context-variable-unfolding rule.
\item \lead{Thesis} zeta kept as a rule, with an explicit conservativity argument for expanding
it away.
\item \lead{lean4lean} \texttt{VExpr} (\ref{ind:vexpr}) has six constructors and no let; the
translation extends the context and returns the body's translation.
\item \lead{Kind} DESIGN.
\item \lead{Why} the thesis's own conservativity argument licenses moving the expansion to the
syntax, keeping the model at six constructors.
\item \lead{Consequence} the kernel's zeta step must be proved to preserve \emph{translation}
rather than to be a rule of the judgment.
\end{itemize}
\paragraph{K-A8: projections are recursor applications, not a term former}
\begin{itemize}
\item \lead{PCUIC} \texttt{tProj} with its own typing rule, its own iota reduction, and its own
congruence rule.
\item \lead{Thesis} a Sigma-projection typing shape, with an eta rule the text marks as not a
real Lean definitional equality.
\item \lead{lean4lean} \texttt{TrProj} (\ref{def:trproj}, \ref{struct:trprojctor}) builds an
application of the structure's own recursor, chaining earlier fields' projections into later
motives.
\item \lead{Kind} DESIGN.
\item \lead{Why} reuses the iota registry (K-A10) and the existing metatheory instead of adding
a node with its own typing, iota and eta rules.
\item \lead{Consequence} \texttt{TrProj} carries \texttt{HasType} side conditions inside a
translation relation, unlike every other clause; coverage is partial (K-A9).
\end{itemize}
\paragraph{K-A9: the projection expansion covers only some structures}
\begin{itemize}
\item \lead{PCUIC} every declared projection typechecks, with no restriction on the structure's
shape.
\item \lead{Thesis} no counterpart.
\item \lead{lean4lean} \texttt{inferProj.WF\_struct} (\ref{thm:vtc-inferproj-struct}) is stated
only for a non-mutual, single-constructor, non-indexed structure, and is \texttt{sorry}; the
general \texttt{inferProj.WF} (\ref{thm:vtc-inferproj}) is not provable as stated.
\item \lead{Kind} GAP.
\item \lead{Why} the recursor-expansion of K-A8 exists only where the generated recursor has the
assumed shape; reflexive, indexed or nested structures have none.
\item \lead{Consequence} a kernel-accepted term projecting such a structure has no translation
derivation at all, so no statement about it can be formed.
\end{itemize}
\paragraph{K-A10: iota as a schematic registry; K-like reduction absent}
\begin{itemize}
\item \lead{PCUIC} one \texttt{cumul\_iota} rule, plus a fix rule guarded by
\texttt{is\_constructor}; no K.
\item \lead{Thesis} the constructor rule, plus K-like reduction for subsingleton eliminators,
tied to axiom K.
\item \lead{lean4lean} \texttt{VEnv} carries a \texttt{pats} registry (\ref{rule:isdefeq-pat});
\texttt{addRecRule} registers only the constructor rule.
\item \lead{Kind} DESIGN for the registry, GAP for K-like reduction.
\item \lead{Why} one generic rule with a registry gives every recursor a reduction without a
per-block case; K-like reduction needs a proof-irrelevance reconstruction the registry cannot
express.
\item \lead{Consequence} the kernel reduces \texttt{Eq.rec} on a non-constructor proof, the
model does not: \texttt{reduceRecursor.WF} (\ref{thm:vtc-reducerecursor}) is \texttt{sorry}.
\end{itemize}
\paragraph{K-A11: quotients are a delta axiom, not an iota rule}
\begin{itemize}
\item \lead{PCUIC} no quotient types.
\item \lead{Thesis} \texttt{Quot}, \texttt{mk}, \texttt{sound}, \texttt{lift}, with only
\texttt{sound} a real axiom -- the rest "semi-builtin".
\item \lead{lean4lean} \texttt{quotDefEq} is one \texttt{VDefEq} entry, fired by
\texttt{IsDefEq.extra}, not a \texttt{pats} entry (\ref{thm:addquot-wf}).
\item \lead{Kind} DESIGN.
\item \lead{Why} a single closed equation with no pattern variables needs no \texttt{Pattern}/
\texttt{PatWF} machinery.
\item \lead{Consequence} the confluence apparatus of K-B6 needs every axiom realised by a
\texttt{pats} rule, which fails as soon as an environment declares the quotient rule.
\end{itemize}
\paragraph{K-A12: inductive-block well-formedness is asserted, not derived}
\begin{itemize}
\item \lead{PCUIC} positivity, universe and elimination conditions, established for a real
environment by \texttt{check\_wf\_env}.
\item \lead{Thesis} generated from the block's specification, not pinned as fixed fields.
\item \lead{lean4lean} \texttt{VInductDecl.WF} (\ref{structure:ind-wf}) is a 20-field syntactic
specification; nothing derives it from the kernel's own inductive checker.
\item \lead{Kind} GAP.
\item \lead{Why} deriving it from \texttt{addInductive} is the missing \texttt{inductDecl} case
of environment admission (K-E3).
\item \lead{Consequence} every theorem about inductive blocks is conditional on a hypothesis no
proof supplies for a kernel-built environment.
\end{itemize}
\paragraph{K-A13: no literal node; primitives verified as reflection}
\begin{itemize}
\item \lead{PCUIC} \texttt{tPrim} as a term former, with its own typing rule and invariants.
\item \lead{Thesis} no counterpart.
\item \lead{lean4lean} no literal constructor in \texttt{VExpr}; a literal unfolds to
constructors, and accelerated arithmetic closes into \texttt{Reflects} statements
(\ref{vpc:thm:reflects-toconst}).
\item \lead{Kind} DESIGN.
\item \lead{Why} keeps the model at six constructors; acceleration is a property of a
declaration the environment already holds.
\item \lead{Consequence} eighteen per-primitive obligations appear with no counterpart on either
reference side, all tainted through the environment's strengthened invariant.
\end{itemize}
\paragraph{K-A14: native reduction refused}
\begin{itemize}
\item \lead{PCUIC} no VM or native conversion modelled.
\item \lead{Thesis} no counterpart.
\item \lead{lean4lean} \texttt{reduceNative} does not support \texttt{reduceBool}; the
obligation (\ref{thm:vtc-reducenative}) is vacuous, the branch never fires.
\item \lead{Kind} OPEN.
\item \lead{Why} verified compilation of native operations would be a project-sized addition on
its own.
\item \lead{Consequence} lean4lean rejects declarations the real kernel accepts through native
reduction; soundness is unaffected, only coverage is reduced.
\end{itemize}

\subsection{K-B: judgments}

\par\smallskip\noindent
\begin{tabular}{lp{0.12\linewidth}p{0.31\linewidth}p{0.27\linewidth}l}
\toprule
id & notion & PCUIC/thesis & lean4lean & kind \\
\midrule
K-B1 & judgment shape & typing and \texttt{cumulSpec0} joined by one rule; thesis keeps them
separate too & one relation \texttt{IsDefEq}, typing its diagonal & DESIGN \\
K-B2 & typed conversion & \texttt{cumulSpec0} is untyped & \texttt{IsDefEq} carries the type;
\texttt{IsDefEqU} forgets it & FORCED-LEAN \\
K-B3 & binder names & \texttt{aname} on every binder, checked in every congruence & \texttt{VExpr}
has no names at all & DESIGN \\
K-B4 & algorithmic equality & a full algorithmic spec, proved equivalent; thesis:
$\Leftrightarrow$ non-transitive, fails SR & no $\Leftrightarrow$ counterpart; stand-in is the
checker's own chain & GAP \\
K-B5 & bidirectional layer & \texttt{Bidirectional/}, \texttt{infer} stated against it &
\texttt{InferType} exists (\ref{thm:infertype-exists}) but unwired to the executable checker &
GAP \\
K-B6 & reduction relation & \texttt{red1} first-class, rules of \texttt{cumulSpec0}; thesis:
$\kappa$ on rec-normal forms & \texttt{ParRed}/\texttt{CParRed}/\texttt{WHRed}/\texttt{StRed}
under an abstract class & DESIGN \\
K-B7 & environments & one \texttt{wf Sigma} predicate for everything & two notions,
\texttt{WF'} and \texttt{Ordered}, bridged silently & DESIGN \\
\bottomrule
\end{tabular}\par\smallskip
\paragraph{K-B1: one relation instead of two}
\begin{itemize}
\item \lead{PCUIC} \texttt{typing} and \texttt{cumulSpec0} as separate inductives, joined by one
conversion rule.
\item \lead{Thesis} typing and definitional equality likewise kept separate.
\item \lead{lean4lean} one four-place \texttt{IsDefEq}\,$E\,U\,\Gamma\,e_1\,e_2\,A$
(\ref{def:isdefeq}), with typing as its diagonal.
\item \lead{Kind} DESIGN.
\item \lead{Why} proof irrelevance and eta already make the two mutually dependent; merging
turns a mutual induction into one.
\item \lead{Consequence} every structural lemma is proved once; a syntax-directed strong
judgment (\ref{def:isdefeqstrong}) is reconstructed separately where inversion needs it.
\end{itemize}
\paragraph{K-B2: conversion is typed}
\begin{itemize}
\item \lead{PCUIC} \texttt{cumulSpec0} relates two terms with no type index.
\item \lead{Thesis} $\equiv$ likewise carries no type, but $\eta$/irrelevance still read one
implicitly.
\item \lead{lean4lean} \texttt{IsDefEq} carries the type as its last argument
(\ref{def:isdefeq}); the untyped form \texttt{IsDefEqU} is derived, agreeing only via
\texttt{thm:unique-typing}.
\item \lead{Kind} FORCED-LEAN.
\item \lead{Why} \texttt{proofIrrel} and \texttt{eta} read the type of the sides they relate.
\item \lead{Consequence} composing equalities at different types needs a family of lemmas that
depend on unique typing, where \texttt{cumul\_Trans} is unconditional in PCUIC.
\end{itemize}
\paragraph{K-B3: no binder names in the model}
\begin{itemize}
\item \lead{PCUIC} \texttt{aname} on every binder; \texttt{eq\_binder\_annot} in every
congruence rule.
\item \lead{Thesis} named binders throughout, matching Lean's real syntax.
\item \lead{lean4lean} \texttt{VExpr}'s \texttt{lam}/\texttt{forallE} carry no name and no
\texttt{BinderInfo} (\ref{family:expr-beq}).
\item \lead{Kind} DESIGN.
\item \lead{Why} names carry no logical content; dropping them removes a premise from three
congruence rules.
\item \lead{Consequence} the kernel's own equality on names/\texttt{BinderInfo} needs a separate
\texttt{EquivBEq} lawfulness argument, not a full \texttt{LawfulBEq}.
\end{itemize}
\paragraph{K-B4: algorithmic equality is not formalised}
\begin{itemize}
\item \lead{PCUIC} a full algorithmic specification, proved equivalent to the declarative one.
\item \lead{Thesis} $\Leftrightarrow$: no transitivity, a head-reduction escape hatch, proved
non-transitive and shown to fail subject reduction.
\item \lead{lean4lean} no relation plays $\Leftrightarrow$'s role; its stand-in is the checker's
own correctness chain.
\item \lead{Kind} GAP.
\item \lead{Why} with $\Leftrightarrow$ unformalised there is no object to state the thesis's
negative results about.
\item \lead{Consequence} nothing marks that a checker-accepted pair need only be algorithmically,
not definitionally, equal.
\end{itemize}
\paragraph{K-B5: the bidirectional layer exists but is unwired}
\begin{itemize}
\item \lead{PCUIC} \texttt{Bidirectional/} proves an inference judgment equivalent to
declarative typing; \texttt{infer} is stated against it.
\item \lead{Thesis} no separate algorithmic layer of this kind.
\item \lead{lean4lean} \texttt{VEnv.InferType} (\ref{thm:infertype-exists}) is sound, complete
and deterministic, but the executable checker is verified against \texttt{IsDefEq} directly.
\item \lead{Kind} GAP.
\item \lead{Why} \texttt{InferType} lives under the abstract reduction class of K-B6, which no
realistic environment satisfies.
\item \lead{Consequence} the completeness route is present in outline and unusable in practice;
\texttt{InferType.exists} is itself doubly tainted.
\end{itemize}
\paragraph{K-B6: reduction is abstract and parameterised}
\begin{itemize}
\item \lead{PCUIC} \texttt{red1} is first-class; \texttt{cumulSpec0} carries the reductions
directly as rules, unconditionally over \texttt{wf Sigma}.
\item \lead{Thesis} $\kappa$ restricts reduction to $\beta/\delta/\zeta/\iota$ on rec-normal-form
terms.
\item \lead{lean4lean} \texttt{ParRed}/\texttt{CParRed}/\texttt{WHRed}/\texttt{StRed} all live
under one abstract \texttt{Params} class fixing an environment and a rule family.
\item \lead{Kind} DESIGN.
\item \lead{Why} abstracting over the rule family proves confluence once for delta and iota
alike.
\item \lead{Consequence} confluence is only as good as the one instance, which requires every
axiom to be realised by a \texttt{pats} rule (K-C1).
\end{itemize}
\paragraph{K-B7: two environment invariants, bridged silently}
\begin{itemize}
\item \lead{PCUIC} \texttt{wf Sigma} is one predicate, \texttt{on\_global\_env}, carrying
everything.
\item \lead{Thesis} a single well-formedness notion, matching PCUIC's.
\item \lead{lean4lean} \texttt{VEnv.WF'} (build-in-order) and \texttt{Ordered} (the weaker
invariant structural lemmas consume), bridged by a coercion instance
(\ref{instance:coeout-orderedstrong}).
\item \lead{Kind} DESIGN.
\item \lead{Why} \texttt{Ordered} is cheap and monotone, so weakening and substitution avoid
carrying declaration history.
\item \lead{Consequence} the coercion hides a \texttt{sorry}-backed hypothesis at every call
site with no visible cue.
\end{itemize}

\subsection{K-C: metatheory}

\par\smallskip\noindent
\begin{tabular}{lp{0.15\linewidth}p{0.34\linewidth}p{0.26\linewidth}l}
\toprule
id & notion & PCUIC/thesis & lean4lean & kind \\
\midrule
K-C1 & confluence & unconditional, triangle method; thesis: Church-Rosser for $\kappa$ up to
$\equiv_p$ & \texttt{church\_rosser} under \texttt{Params}, two open sorries & GAP \\
K-C2 & subject reduction & \texttt{subject\_reduction}, proved; thesis calls it "an easy
induction" & built into the \texttt{pat} rule; the honest form is \texttt{sorry} & GAP \\
K-C3 & principality & proved without normalisation; thesis: unique typing via
$\vdash_n$-stratification & \texttt{IsDefEq.uniq} via \texttt{HasTypeStratified}, tainted twice &
GAP \\
K-C4 & type-former injectivity & derived from confluence, proved; thesis's \texttt{thm:1dinv}
closes its induction & all three injectivity principles are \texttt{sorry} & GAP \\
K-C5 & strong normalisation & a typeclass axiom, used once; thesis: an unfinished sketch, no
theorem & none; the checker runs on explicit fuel & DESIGN \\
K-C6 & canonicity, consistency & proved, under normalisation; thesis: ZFC + inaccessibles, on
paper & none; only soundness relative to the model & GAP \\
K-C7 & evaluation & weak-CBV evaluation, \texttt{wcbv\_standardization}; thesis has no
standardization theorem & \texttt{ParRedS.standard}, about the ideal reduction only & GAP \\
\bottomrule
\end{tabular}\par\smallskip
\paragraph{K-C1: confluence holds only for an unrealistic class of environments}
\begin{itemize}
\item \lead{PCUIC} the triangle method with an optimal-reduction function, proved
unconditionally for any well-formed environment.
\item \lead{Thesis} the analogue for $\kappa$, up to $\equiv_p$, via the Tait--Martin-L\"of
method.
\item \lead{lean4lean} \texttt{ParRedS.church\_rosser} holds up to \texttt{NormalEq}, but under
\texttt{Params} and with two literal sorries in a supporting lemma.
\item \lead{Kind} GAP.
\item \lead{Why} the \texttt{Params} restriction is the price of K-B6's abstraction; the two
sorries predate the branches.
\item \lead{Consequence} Church-Rosser, standardisation, and the head-reduction lemmas are all
unproved for any environment with a definition or the quotient rule.
\end{itemize}
\paragraph{K-C2: subject reduction for iota is a rule, not a theorem}
\begin{itemize}
\item \lead{PCUIC} \texttt{subject\_reduction} and its one-step form \texttt{sr\_red1}, both
proved.
\item \lead{Thesis} the corresponding fact, dismissed as "an easy induction."
\item \lead{lean4lean} \texttt{IsDefEq.pat} asserts the reduct's type with no typing premise on
it; the honest form needs \texttt{VEnv.WF.patsStrong} (\ref{thm:wf-patsstrong}), which is
\texttt{sorry}.
\item \lead{Kind} GAP (for the obligation), DESIGN (for the rule's shape).
\item \lead{Why} building it into the rule lets \texttt{IsDefEq} exist before the strong system
does; the repayment is deferred.
\item \lead{Consequence} every declaration depending on \texttt{VEnv.WF} becomes conditional on
this one obligation.
\end{itemize}
\paragraph{K-C3: unique typing by a different route, and tainted}
\begin{itemize}
\item \lead{PCUIC} \texttt{principal\_type}, proved explicitly without relying on normalisation.
\item \lead{Thesis} \texttt{thm:unique} via a $\vdash_n$-indexed joint induction with
Church-Rosser.
\item \lead{lean4lean} \texttt{IsDefEq.uniq} (\ref{thm:unique-typing}) strengthens, projects to a
syntax-directed judgment, then stratifies by derivation level.
\item \lead{Kind} GAP (for the taint), DESIGN (for the route).
\item \lead{Why} stratifying by derivation level avoids the thesis's mutual recursion between
typing and conversion.
\item \lead{Consequence} unique typing is available conditionally; because K-B2 routes
compositions of equalities through it, the taint is pervasive.
\end{itemize}
\paragraph{K-C4: type-former injectivity is entirely open}
\begin{itemize}
\item \lead{PCUIC} sort and product injectivity under conversion, derived from confluence and
proved.
\item \lead{Thesis} \texttt{thm:1dinv}, the step closing the unique-typing induction.
\item \lead{lean4lean} three declarations (\ref{thm:injectivity-open}), all \texttt{sorry}; the
derived \texttt{forallE\_inv} inherits both gaps.
\item \lead{Kind} GAP.
\item \lead{Why} two routes are in flight -- the blocked Church-Rosser line, and an experimental
logical-relations line with its own \texttt{sorry}.
\item \lead{Consequence} unique typing, and the head-reduction lemmas that expose a type's head
shape, all appeal to an already-open foundation.
\end{itemize}
\paragraph{K-C5: no normalization statement; the checker runs on fuel}
\begin{itemize}
\item \lead{PCUIC} normalisation as two typeclass assumptions, used to justify a fuel-free
weak-head reduction.
\item \lead{Thesis} an unfinished sketch of a terminating variant, stopping before any theorem.
\item \lead{lean4lean} several independent fuel counters (\ref{thm:vtc-withfuel}); a
fuel-exhausted run is unconstrained.
\item \lead{Kind} DESIGN.
\item \lead{Why} Lean definitions need explicit termination witnesses, and a fuel counter is one.
\item \lead{Consequence} the largest single axiom of the Rocq boundary is absent, but
completeness becomes unstatable in the current shape.
\end{itemize}
\paragraph{K-C6: no canonicity, no consistency}
\begin{itemize}
\item \lead{PCUIC} canonicity and consistency, both proved under the normalisation assumption.
\item \lead{Thesis} soundness of a proof-split variant in ZFC plus countably many
inaccessibles, on paper, for a slightly stronger system.
\item \lead{lean4lean} none; the project states correctness relative to the model, not
consistency of the type theory.
\item \lead{Kind} GAP.
\item \lead{Why} out of scope by design: the statement is that the checker accepts a term only
if the model derives its typing.
\item \lead{Consequence} the trust chain terminates at \texttt{VEnv.IsDefEq}; nothing states
that relation is consistent.
\end{itemize}
\paragraph{K-C7: standardization is about the ideal reduction only}
\begin{itemize}
\item \lead{PCUIC} one-shot big-step weak call-by-value evaluation, with
\texttt{wcbv\_standardization} the sole use of strong normalisation in this part of the
development.
\item \lead{Thesis} no standardization theorem at all.
\item \lead{lean4lean} \texttt{ParRedS.standard} (\ref{thm:parreds-standard}), after Kashima
(2000), about the abstract ideal reduction, under \texttt{Params}.
\item \lead{Kind} GAP.
\item \lead{Why} the theorem serves the head-reduction lemmas exposing a type's head shape, not
an observational statement about running programs.
\item \lead{Consequence} no weak call-by-value evaluation relation, and no value predicate,
exists for Lean terms anywhere in the development.
\end{itemize}

\subsection{K-D: checker and method}

\par\smallskip\noindent
\begin{tabular}{lp{0.12\linewidth}p{0.27\linewidth}p{0.25\linewidth}l}
\toprule
id & notion & PCUIC/thesis & lean4lean & kind \\
\midrule
K-D1 & checker guarantee & sound and complete; thesis proves $\equiv$ undecidable,
$\Leftrightarrow$ non-transitive & sound only; failure of any kind is unconstrained & FORCED-LEAN
\\
K-D2 & verification method & checker equals spec, same terms, extracted; no thesis analogue &
a port of the C++ kernel, verified through a refinement layer & DESIGN \\
K-D3 & open obligations & none on the Rocq side; no thesis analogue & six sorries in
\texttt{Verify/TypeChecker/} & GAP \\
K-D4 & pointer equality & none, a pure function; no thesis analogue & one axiom the algorithm's
fast path branches on & FORCED-LEAN4LEAN \\
K-D5 & host data structures & defined in Rocq, nothing assumed; no thesis analogue & 32 bridge
axioms over Lean's own unverified containers & FORCED-LEAN4LEAN \\
K-D6 & guard condition & two named axioms, the largest single hole; thesis: recursion compiled
before the kernel & nothing to axiomatise; no fix constructor at all & FORCED-LEAN \\
K-D7 & termination & fuel-free, well-founded on strong normalisation; no thesis analogue &
several independent fuel counters & DESIGN \\
K-D8 & level decision procedure & graph acyclicity, one algorithm; thesis: a rule system, no
algorithm stated & a canonical form, proved complete, plus a coNP-hardness result & DESIGN \\
\bottomrule
\end{tabular}\par\smallskip
\paragraph{K-D1: soundness only, and completeness is not merely unproved}
\begin{itemize}
\item \lead{PCUIC} sound \emph{and} complete against the declarative specification, every
rejection carrying a refutation.
\item \lead{Thesis} proves $\equiv$ undecidable and $\Leftrightarrow$ non-transitive, so no
algorithm decides either.
\item \lead{lean4lean} the three public entry points (\ref{thm:vtc-toplevel}) give soundness
only: a \texttt{true} result relates the two terms, nothing else is constrained.
\item \lead{Kind} FORCED-LEAN.
\item \lead{Why} any completeness claim would have to be relative to $\Leftrightarrow$, which
K-B4 says is not formalised.
\item \lead{Consequence} a rejection carries no information; nothing detects an incompleteness
bug of the kind the Rocq checker's method found.
\end{itemize}
\paragraph{K-D2: verification by translation to a model, not by construction}
\begin{itemize}
\item \lead{PCUIC} the checker is written in Rocq against the same syntax its specification
uses; checked and specified object coincide.
\item \lead{Thesis} no checker-construction analogue; describes only the ideal judgments.
\item \lead{lean4lean} a port of Lean's C++ kernel on the real \texttt{Lean.Expr}; correctness
goes through \texttt{TrExprS}/\texttt{TrEnv} and a companion \texttt{.WF} lemma per function.
\item \lead{Kind} DESIGN.
\item \lead{Why} fidelity: the verified object is the algorithm that actually runs, at the
performance of the real kernel.
\item \lead{Consequence} an entire refinement layer exists with no counterpart, bringing its own
trust boundary (K-D4, K-D5) and its own vacuity risks (K-E3).
\end{itemize}
\paragraph{K-D3: six open obligations in the checker's verification}
\begin{itemize}
\item \lead{PCUIC} no open per-function obligations.
\item \lead{Thesis} no checker-construction analogue.
\item \lead{lean4lean} six sorries: \ref{thm:vtc-tryetastruct}, \ref{thm:vtc-isdefequnitlike}
(K-A6); \texttt{reduceProjCore.WF}; \ref{thm:vtc-reducerecursor} (K-A10);
\ref{thm:vtc-inferproj-struct}, \ref{thm:vtc-inferproj} (K-A9).
\item \lead{Kind} GAP.
\item \lead{Why} four of the six trace to K-A6 and K-A8/K-A9: the model has no structure eta and
no projection node.
\item \lead{Consequence} every statement in the top-level triple is partial as stated; the iota
step itself is proved but not wired into the recursor obligation.
\end{itemize}
\paragraph{K-D4: pointer equality as an axiom}
\begin{itemize}
\item \lead{PCUIC} the checker is a pure Rocq function; no physical-identity test anywhere.
\item \lead{Thesis} no checker-construction analogue.
\item \lead{lean4lean} \texttt{ptrEqExpr\_eq} axiomatises that pointer-equal expressions are
equal; the \texttt{EquivManager} cache built on top is verified separately
(\ref{thm:eqvmgr-isdefeq-wf}).
\item \lead{Kind} FORCED-LEAN4LEAN.
\item \lead{Why} the ported algorithm branches on pointer identity for performance; the axiom
states what that branch assumes.
\item \lead{Consequence} mostly benign, except one call site where the axiom gates a different
algorithm, not merely a shortcut.
\end{itemize}
\paragraph{K-D5: bridge axioms over the host's own data structures}
\begin{itemize}
\item \lead{PCUIC} environments, maps and arrays are all defined in Rocq; nothing is assumed
about them.
\item \lead{Thesis} no checker-construction analogue.
\item \lead{lean4lean} 32 axioms relate the model to Lean's own \texttt{PersistentHashMap},
\texttt{PersistentArray}, \texttt{Expr}, \texttt{Level} and \texttt{Syntax}.
\item \lead{Kind} FORCED-LEAN4LEAN.
\item \lead{Why} the kernel runs on Lean's own container implementations; verifying those is a
separate project.
\item \lead{Consequence} the trust boundary is a mixture of mathematical gaps and library
lemmas, not a short list of axioms as on the Rocq side.
\end{itemize}
\paragraph{K-D6: no guard condition to axiomatise}
\begin{itemize}
\item \lead{PCUIC} \texttt{guard\_checking} and \texttt{guard\_checking\_correct}, its largest
single trust axiom.
\item \lead{Thesis} recursion is compiled into recursor applications before the kernel sees a
declaration.
\item \lead{lean4lean} nothing: \texttt{VExpr} (\ref{ind:vexpr}) has no fix constructor and
\texttt{IsDefEq} has no fix rule.
\item \lead{Kind} FORCED-LEAN.
\item \lead{Why} Lean's design puts termination checking outside the kernel entirely.
\item \lead{Consequence} the Lean side's named-axiom boundary is smaller here; with K-C5 removing
the other Rocq axiom too, neither has a direct counterpart.
\end{itemize}
\paragraph{K-D7: fuel instead of well-founded recursion on normalisation}
\begin{itemize}
\item \lead{PCUIC} \texttt{reduce\_stack} is fuel-free, well-founded on an order whose
accessibility comes from the normalisation axiom.
\item \lead{Thesis} no checker-construction analogue.
\item \lead{lean4lean} several independent fuel counters (\ref{thm:vtc-withfuel}); the real
kernel uses one shared recursion-depth bound.
\item \lead{Kind} DESIGN.
\item \lead{Why} Lean definitions need explicit termination witnesses, and a fuel counter is
one.
\item \lead{Consequence} the two implementations can disagree on deeply recursive input purely
by bound, and exhaustion is indistinguishable from rejection.
\end{itemize}
\paragraph{K-D8: a complete level decision procedure, and a hardness result}
\begin{itemize}
\item \lead{PCUIC} weighted-graph acyclicity, deciding constraint consistency, not equivalence
of level expressions.
\item \lead{Thesis} the rule system of K-A2, with no algorithm and no complexity result stated.
\item \lead{lean4lean} a canonical form proved sound and complete (\ref{thm:lvl-normalize-complete}),
plus a coNP-hardness result (\ref{thm:levelsat-equiv-iff-unsat}) for the decision problem.
\item \lead{Kind} DESIGN.
\item \lead{Why} the semantic order of K-A2 must be decided completely, not merely
approximated.
\item \lead{Consequence} lean4lean decides level order for strictly more pairs than the real
kernel, and never the converse, a deliberate one-directional divergence.
\end{itemize}

\subsection{K-E: the trust boundary}

\par\smallskip\noindent
\begin{tabular}{llp{0.27\linewidth}p{0.32\linewidth}l}
\toprule
id & notion & PCUIC/thesis & lean4lean & kind \\
\midrule
K-E1 & boundary shape & two named axioms plus quoting; no thesis analogue & 16 live sorries,
110 axioms, 12 opaque, quoted per statement & GAP \\
K-E2 & reification & Template quoting, unverified; no thesis analogue & none; \texttt{Replay}
reads the real environment directly & DESIGN \\
K-E3 & vacuous layers & none reported; no thesis analogue & three proved apparatuses, each
inhabited by nothing reachable & GAP \\
\bottomrule
\end{tabular}\par\smallskip
\paragraph{K-E1: the shape of the trust boundary}
\begin{itemize}
\item \lead{PCUIC} the guard condition and normalisation, two named axioms, plus the unverified
quoting layer and OCaml extraction.
\item \lead{Thesis} no boundary statement of this kind; describes the ideal judgments only.
\item \lead{lean4lean} 16 live sorries outside \texttt{Experimental/}, 110 axioms and 12 opaque
constants, quoted per top-level statement.
\item \lead{Kind} GAP.
\item \lead{Why} the refinement method of K-D2 spreads the boundary across theory gaps and
library assumptions, where a correct-by-construction checker has only the first.
\item \lead{Consequence} a reader cannot state what lean4lean assumes in one line; roughly 8\%
of user-level declarations are \texttt{sorryAx}-tainted.
\end{itemize}
\paragraph{K-E2: no reification layer}
\begin{itemize}
\item \lead{PCUIC} the Template layer quotes Rocq's own term representation into the checked
syntax, in OCaml, unverified.
\item \lead{Thesis} no reification step; the ideal judgments are stated directly on Lean-shaped
terms.
\item \lead{lean4lean} \texttt{Replay.lean} reads declarations out of a real imported
environment and replays them through the executable kernel.
\item \lead{Kind} DESIGN.
\item \lead{Why} Lean's syntax is already a Lean datatype; no reification is needed.
\item \lead{Consequence} one item of the Rocq boundary closes outright, at the cost of depending
on Lean's own container implementations (K-D5).
\end{itemize}
\paragraph{K-E3: three vacuous layers}
\begin{itemize}
\item \lead{PCUIC} no layer is proved yet inhabited by nothing; theorems are stated over a
predicate a real checker establishes.
\item \lead{Thesis} no counterpart.
\item \lead{lean4lean} the \texttt{inductDecl} case of environment admission is \texttt{sorry}
(\ref{structure:ind-wf}); the confluence class needs every axiom realised as a pattern
(\ref{thm:indparams-defeqs-as-pats-vacuous}); \texttt{TrEnv.proj\_defeq}
(\ref{thm:tr-env-proj-defeq}) has no caller.
\item \lead{Kind} GAP.
\item \lead{Why} in each case the specification side is finished and the refinement side that
would inhabit it is not.
\item \lead{Consequence} counting proved declarations overstates what is established; a vacuity
pass is the only honest reading of the census.
\end{itemize}

\section{Thesis versus code}
\label{sec:ref-thesis-code}

Thirteen further items compare Carneiro's thesis directly against the code, where the code
either departs from the thesis by design or goes beyond it; each is cross-referenced to the
K-A--K-E item that treats it in full.

\par\smallskip\noindent
\begin{tabular}{lp{0.12\linewidth}p{0.20\linewidth}p{0.19\linewidth}lp{0.22\linewidth}}
\toprule
id & notion & thesis & code & kind & why \\
\midrule
K-T1 & levels & the nine-rule system of \texttt{axioms.tex} & \texttt{VLevel.LE}, a semantic
order & DESIGN & no rule-completeness proof needed (K-A2) \\
K-T2 & level complexity & no complexity result stated & \texttt{LevelSat.lean} proves the
problem coNP-hard & DESIGN & bounds what any complete algorithm can cost (K-A2, K-D8) \\
K-T3 & let/zeta & zeta kept, with a conservativity argument & no zeta rule; let substituted away
& DESIGN & keeps the model at six constructors (K-A7) \\
K-T4 & algorithmic equality & $\Leftrightarrow$: undecidable, non-transitive, fails subject
reduction & $\Leftrightarrow$ has no formal counterpart at all & GAP & the thesis's negative
results have no formal home (K-B4, K-D1) \\
K-T5 & unique typing & $\vdash_n$-indexed stratification & \texttt{HasTypeStratified}, a strong
syntax-directed judgment & DESIGN & replaces alternation-counting with level-counting (K-C3) \\
K-T6 & rec-normal forms & the eta-vs-iota fight on partially applied eliminators & no
counterpart & GAP & the fight is not modelled at all \\
K-T7 & eight primitives & \texttt{Wtypes.tex}'s $\Sigma$/$W$/\texttt{acc} encoding & not
formalised; only the Sigma-projection shape reused & DESIGN & the shape serves \texttt{TrProj},
the reduction does not (K-A8) \\
K-T8 & soundness & a ZFC-plus-inaccessibles model, on paper & no counterpart & GAP & correctness
is stated relative to the model, not consistency (K-C6) \\
K-T9 & iota vs.\ K-split & \texttt{normalization.tex} splits $K^+$ from $\iota$ & one uniform
iota rule & DESIGN & matches the stub's weaker relation, not its split (K-A10) \\
K-T10 & compilation stub & a 25-line fragment, no theorem & native reduction refused outright &
OPEN & no verified-compilation route exists yet (K-A14) \\
K-T11 & inductive blocks & the recursor and iota rules generated from the block's spec &
\texttt{VInductDecl.WF} pins 20 fields directly & DESIGN & makes well-formedness decidable, at
the cost of K-A12 \\
K-T12 & standardization & no theorem stated & \texttt{ParRedS.standard}, after Kashima (2000) &
DESIGN & beyond the thesis; serves the head-reduction lemmas (K-C7) \\
K-T13 & regularity of reductions & called "an easy induction" & \texttt{VEnv.WF.patsStrong}, the
one open obligation & GAP & the widest-reaching gap in the project (K-C2) \\
\bottomrule
\end{tabular}\par\smallskip

\section{What is aligned}
\label{sec:ref-aligned}

The remaining notions line up name for name, with no divergence of kind; each row is a MetaRocq
or thesis element and its lean4lean counterpart.

\par\smallskip\noindent
\begin{tabular}{p{0.23\linewidth}p{0.19\linewidth}p{0.27\linewidth}p{0.20\linewidth}}
\toprule
notion & PCUIC & lean4lean & thesis \\
\midrule
variable lookup & \texttt{type\_Rel} & \texttt{Lookup}, \texttt{IsDefEq.bvar} & the variable rule
\\
sorts & \texttt{type\_Sort} & \texttt{IsDefEq.sortDF} & -- \\
product formation & \texttt{type\_Prod}, \texttt{sort\_of\_product} & \texttt{IsDefEq.forallEDF}
with \texttt{imax} & Pi-formation, the \texttt{imax} clause \\
abstraction & \texttt{type\_Lambda} & \texttt{IsDefEq.lamDF} & -- \\
application & \texttt{type\_App} & \texttt{IsDefEq.appDF} & -- \\
universe-polymorphic constants & \texttt{type\_Const} & \texttt{IsDefEq.constDF} & -- \\
conversion & \texttt{type\_Cumul} & \texttt{IsDefEq.defeqDF} (modulo K-A3) & -- \\
context well-formedness & \texttt{wf\_local} & \texttt{OnCtx}\,$\Gamma\,(\mathrm{IsType})$ &
$\vdash \Gamma$ ok \\
beta & \texttt{cumul\_beta} & \texttt{IsDefEq.beta} & -- \\
delta & \texttt{cumul\_delta} & \texttt{IsDefEq.extra} over \texttt{env.defeqs} & definitions \\
iota & \texttt{cumul\_iota} & \texttt{IsDefEq.pat} over \texttt{env.pats} (modulo K-A10) & the
iota section \\
weakening & \texttt{weakening\_typing} & \texttt{IsDefEq.weakN} & \texttt{thm:weak} \\
substitution & \texttt{substitution} & \texttt{IsDefEq.instN} & \texttt{thm:subst}, item
\texttt{subst\_ty} \\
validity, regularity of types & \texttt{validity} & \texttt{IsDefEq.isType} & \texttt{thm:reg},
item 4 \\
free variables in context & \texttt{PCUICClosed} & \texttt{IsDefEq.closedN} & \texttt{thm:reg},
item 2 \\
context conversion & \texttt{PCUICContextConversion} & \texttt{IsDefEqCtx},
\texttt{IsDefEq.defeqDFC} & -- \\
environment weakening & \texttt{weakening\_env} & \texttt{IsDefEq.mono} along the extension order
& -- \\
universe instantiation & \texttt{typing\_subst\_instance} & \texttt{IsDefEq.instL} &
universe-polymorphic constants \\
simultaneous substitution & \texttt{subslet} & \texttt{Ctx.SubstEq}, \texttt{IsDefEq.substDF} &
-- \\
parallel reduction & \texttt{pred1} & \texttt{ParRed} (\ref{def:parred}) & $\rightsquigarrow_\kappa$
(compatibility direction) \\
complete parallel reduction & \texttt{rho} & \texttt{CParRed} (\ref{def:cparred}) & -- \\
triangle lemma & the triangle property & \texttt{ParRed.triangle} & \texttt{thm:tri} \\
confluence & \texttt{PCUICConfluence.v} & \texttt{ParRedS.church\_rosser} (status differs, K-C1)
& \texttt{thm:church\_rosser} \\
defeq via reduction plus irrelevance & -- & \texttt{IsDefEq.church\_rosser}
(\ref{thm:isdefeq-church-rosser}) & \texttt{thm:ckappa} \\
bidirectional inference & \texttt{Bidirectional/} & \texttt{VEnv.InferType} (unwired, K-B5) & --
\\
head-reduce to a sort or product & \texttt{reduce\_to\_sort}/\texttt{reduce\_to\_prod} &
\texttt{IsDefEq.reduce\_sort}/\texttt{.reduce\_forallE} & -- \\
level valuation semantics & \texttt{Universes.val}/\texttt{satisfies} &
\texttt{VLevel.eval}/\texttt{VLevel.LE} (\ref{ind:vlevel}) & level semantics paragraph \\
declaration admission & \texttt{check\_wf\_decl} & \texttt{VDecl.WF}, \texttt{VEnv.WF'} & \S2.5,
\S2.6, \S2.7.1 \\
mutual blocks & MetaRocq's mutual definitions & \texttt{VEnv.addConsts} then
\texttt{addDefEqs} & -- \\
monotonicity of well-formedness & \texttt{weakening\_env\_decl} & \texttt{VConstant.WF.mono},
\texttt{VDefEq.WF.mono}, \texttt{PatWF.mono} & -- \\
\bottomrule
\end{tabular}\par\smallskip

\section*{Caveats}

\begin{itemize}
\item Every claim above is read off a source file, a blueprint label, or the analysis behind
this chapter; none is inferred from a paper's abstract or from a title alone.
\item [CCC] is superseded by [JACM] on every shared topic; it is cited here only for its earlier
framing (PCUIC's introduction) or for historical firsts (Sec 3, first sound checker).
\item [MR] appears only where lean4lean's own evidence explicitly names it; it establishes no
kernel-verification result and should never be cited for one.
\item A "kind" is read off the current state of the code and the current state of the two
papers; a later proof could move a GAP item into DESIGN without changing what is proved today.
\item Kind counts (FORCED-LEAN 7, FORCED-LEAN4LEAN 2, DESIGN 22, GAP 19, OPEN 2) total all 52
items: K-A--K-E (DESIGN 14, GAP 15, OPEN 1) plus the K-T table of \S\ref{sec:ref-thesis-code}
(DESIGN 8, GAP 4, OPEN 1).
\end{itemize}
